\section{Thesis Contributions} \label{sec:15-introduction-contributions}
\begin{foldable}
  The three restricted settings identified in \S\ref{sec:12-introduction-model} and \S\ref{sec:13-introduction-dataset} are few-shot and data-free black-box model distillation, and text dataset distillation.
  Each imposes a distinct form of information scarcity on the distillation process.
  Structurally, these settings map to two complementary axes of \ac{KD}: the first operates on the model axis (Chapters~\ref{ch:30-research01}--\ref{ch:40-research02}), compressing the learned function under restricted teacher access; the second operates on the data axis (Chapter~\ref{ch:50-research03}), compressing the training signal under discrete-token and coverage constraints.
  Across both axes, distributional fidelity is the binding constraint.
  The thesis addresses all three settings under a single design philosophy: treat distributional fidelity as a first-class objective, derive practical constraints from the deployment setting, and enforce representational coverage of the distilled signal explicitly.
  The three contributions are as follows.
\end{foldable}

\begin{foldable}
  \boldtight{Diversity-Aware Black-Box Few-Shot Knowledge Distillation} (\textbf{\acs{DivBFKD}}).
  In Chapter~\ref{ch:30-research01}, \ac{DivBFKD} addresses the setting in which only a small number of real training images are available and the teacher is accessible only as a black box, returning final predictions rather than internal representations.
  The method trains a \ac{WGAN} to synthesize additional training images, augmenting this process with a high-confidence image scheme that identifies and retains synthetic samples whose teacher-assigned confidence exceeds class-specific adaptive thresholds computed via the $q$-quantile of the score distribution.
  This threshold acts as a quality gate, admitting only sufficiently informative and class-coherent synthetic images into the distillation pipeline, and directly counters the semantic narrowness that arises when few real examples are available.
  We evaluate across seven image classification benchmarks and reach state-of-the-art performance, with a gain of over 17\% relative to standard \ac{KD} on CIFAR-10 and closing to within 2\% of a fully supervised student upper bound on the simpler datasets.
  This work was published at ECML-PKDD 2024.
\end{foldable}

\begin{foldable}
  \boldtight{Diverse Image Priors for Black-Box Data-Free Knowledge Distillation} (\textbf{\acs{DIPKD}}).
  Later, in Chapter~\ref{ch:40-research02}, \ac{DIPKD} addresses the strictest image-domain setting: the student has access to neither real training data nor teacher internals, receiving only top-1 hard labels in response to queries.
  The method is organised as a three-phase pipeline.
  In the first phase, a diverse pool of synthetic images is built using hierarchical noise injection, nonlinear photometric transforms, and semantic cutmixing, producing candidate images that cover a broad region of the input space.
  In the second phase, a lightweight primer student is trained on these candidates under a contrastive objective that maximises inter-class separation, sharpening the diversity of the pool before full distillation begins.
  In the third phase, the trained primer student supplies surrogate logits as soft distillation labels to supplement the teacher's hard labels, recovering much of the information lost to black-box access.
  \ac{DIPKD} reaches state-of-the-art results on 11 of 12 benchmarks, with a gain of over 21\% on Imagenette relative to the best prior baseline, and shows strong performance on medical imaging datasets from MedMNIST.
  This work has been submitted to ICCCI 2026.
\end{foldable}

\begin{foldable}
  \boldtight{Trajectory-Aware Knowledge Estimation for Text Dataset Distillation} (\textbf{\acs{TAKE}}).
  Finally, in Chapter~\ref{ch:50-research03}, \ac{TAKE} addresses dataset distillation in the natural language domain, where the goal is to distill a compact subset of a large text corpus that can surrogates the performance of training on the full dataset.
  Existing text dataset distillation methods fall short on three fronts: their weighting schemes are uniform and uninformative about task-aligned importance, their objectives optimize batch by batch rather than enforcing global distributional coverage, and their outputs are embeddings rather than human-readable text.
  \ac{TAKE} addresses all three through a multi-stage methodology.
  First, it quantifies each sample's downstream contribution via influence functions and integrates these estimates across the full training trajectory.
  Second, it generates a pool of human-readable synthetic candidates through language model fine-tuning, then distills a compact corpus via discrete optimal transport that enforces global distributional coverage rather than relying on per-sample independence.
  Theoretically, \ac{TAKE} formalizes the hard-sample bias in text distillation and derives formal bounds on the downstream loss gap.
  Empirically, \ac{TAKE} matches or exceeds baseline methods across six text datasets at extreme compression ratios, producing human-readable corpora that generalize across diverse model architectures.
  This work has been submitted to ECML-PKDD 2026.
\end{foldable}

\begin{foldable}
  Operationally, the design philosophy becomes concrete mechanisms tailored to the information scarcity each setting imposes.
  \ac{DivBFKD} (image, black-box few-shot) filters synthetic candidates through generation and adaptive confidence thresholds, admitting only semantically coherent, informative examples into distillation.
  \ac{DIPKD} (image, black-box data-free) synthesizes diverse image priors, applies contrastive learning to sharpen distinction between samples, then distills via a primer student that supplies soft signals, addressing the diversity deficit that arises without real data.
  \ac{TAKE} (text dataset) concentrates budget through knowledge-based influence weighting and discrete optimal transport, identifying high-value samples (not uniform, uninformative ones), while trajectory integration corrects hard-sample bias and enforces global coverage.
  Across all three, the unifying insight is the same: representational coverage and quality are not implicit properties of the training signal, but design objectives that require explicit measurement and enforcement.
  Together, these contributions show that \ac{KD} under restricted access is feasible and can match or exceed methods that operate under far less constrained conditions.
\end{foldable}
